\section{A simple module with polynomial self-extensions}
\label{res:section}

The cocycle constructed above will control the two cones in our example.
We now determine all self-extensions of the simple on which it acts.
The key step is a calculation over the smaller algebra $C$: its simple $s$
has no positive self-extensions, but $\Ext_C^2(s,C)$ is a simple module.
Passing to the trivial extension $T=C\ltimes DC$ turns this exceptional
class into a polynomial generator of degree three.

\subsection{The resolution over \texorpdfstring{$C$}{C}}

Let $s^{\mathrm r}$ denote the one-dimensional right $f$-simple.
For every integer $i\geq0$, set
\[
 \ell_i=x+q^i y.
\]
The transcendence of $q$ implies $q\ne0$ and $1+q^m\ne0$ for every
positive integer $m$. We will use these conditions for arbitrarily large
indices.

\begin{lemma}\label{res:base}
A projective resolution $R\to s$ over $C$ has terms
\[
 R_0=Cf,\qquad R_i=Ce\quad(i\geq1),
\]
and differentials given by right multiplication:
\[
 \cdots\longrightarrow Ce
 \xrightarrow{\,\cdot\ell_1\,}Ce
 \xrightarrow{\,\cdot\ell_0\,}Ce
 \xrightarrow{\,\cdot u\,}Cf\longrightarrow s\longrightarrow0.
\]
More precisely, $d_{i+2}$ is right multiplication by $\ell_i$ for every
$i\geq0$. For every integer $i\geq0$,
\[
 \Ext_C^i(s,s)=
 \begin{cases}\kk&i=0,\\0&i>0,\end{cases}
 \qquad
 \Ext_C^i(s,C)=
 \begin{cases}s^{\mathrm r}&i=2,\\0&i\ne2.\end{cases}
\]
The second identification is an identification of right $C$-modules.
\end{lemma}

\begin{proof}
The left ideals $Ce$ and $Cf$ have respective bases
\[
 (e,x,y,z,t,j),\qquad(f,u,v,n).
\]
The multiplication table gives the following formulas, valid for every
$i\geq0$:
\[
\begin{array}{c|cccccc}
 a&e&x&y&z&t&j\\ \hline
 au&u&v&v&0&(1+q)n&0\\
 a\ell_i&\ell_i&q^{i+1}z&z&0&(1+q^{i+2})j&0.
\end{array}
\]
Consequently, right multiplication by $u$ has image
$\langle u,v,n\rangle=\rad(Cf)$ and kernel
$\langle\ell_0,z,j\rangle$. Right multiplication by $\ell_i$ has image
$\langle\ell_i,z,j\rangle$ and kernel
$\langle\ell_{i+1},z,j\rangle$. In the latter calculation, a kernel vector
has zero $e$- and $t$-coefficients, while its $x,y$ coefficients satisfy
$a_y=q^{i+1}a_x$. These identities prove exactness at every index.
Since $\Hom_C(Ce,s)=es=0$ and $\Hom_C(Cf,s)=fs=\kk$, they also give the
asserted self-extension groups over $C$.

To compute $\Ext_C^*(s,C)$, identify
$\Hom_C(Ce,C)=eC$ and $\Hom_C(Cf,C)=fC$ by evaluation at the idempotent.
Precomposition with right multiplication becomes left multiplication.
Thus the Hom complex is
\[
 fC\xrightarrow{\,u\cdot\,}eC
 \xrightarrow{\,\ell_0\cdot\,}eC
 \xrightarrow{\,\ell_1\cdot\,}eC\longrightarrow\cdots.
\]
Left multiplication by $u$ sends the basis $(f,t,j,n)$ of $fC$ to
$(u,y+qx,z,v)$. These four vectors are independent. On the basis
$(e,x,y,z,u,v)$ of $eC$, left multiplication by $\ell_i$ is
\[
\begin{array}{c|cccccc}
 a&e&x&y&z&u&v\\ \hline
 \ell_i a&\ell_i&q^i z&qz&0&(1+q^i)v&0.
\end{array}
\]
The index $i=0$ is exceptional:
\[
 \ker(\ell_0\cdot)=\langle qx+y,z,u,v\rangle=\im(u\cdot),
 \qquad
 \im(\ell_0\cdot)=\langle\ell_0,z\rangle.
\]
For every $i\geq1$, the coefficient $1+q^i$ is nonzero, and
\[
 \ker(\ell_i\cdot)=\langle\ell_{i-1},z,v\rangle,
 \qquad
 \im(\ell_i\cdot)=\langle\ell_i,z,v\rangle.
\]
The cohomology in degrees zero and one therefore vanishes. In degree two
it is the line represented by $v$; in every higher degree the displayed
kernel equals the preceding image. Finally, $vf=v$, $ve=0$, and the only
possibly nonzero right radical product is $vt=qz$, which is a boundary.
The surviving line is precisely $s^{\mathrm r}$.
\end{proof}

\subsection{The trivial-extension triangle and its product}

Write $L=DC$. The duality calculation in Lemma~\ref{res:base} determines
$L\otimes_C^{\mathbf L}s$. The extension $0\to L\to T\to C\to0$
then provides a connecting morphism over $T$. We identify its products
explicitly so that the calculation controls every positive degree.

\begin{lemma}\label{res:polynomial}
For the algebra $T=C\ltimes DC$ and its left $f$-simple $s$,
\[
 \Ext_T^*(s,s)=\kk[\tau],\qquad |\tau|=3,
\]
as a graded algebra, where $\tau=p(s)$ is the evaluation of the cocycle
from Lemma~\ref{coc:data}.
\end{lemma}

\begin{proof}
For a finite projective left $C$-module $V$, evaluation gives a natural
isomorphism of left $C$-modules
\[
 DC\otimes_C V\longrightarrow D\Hom_C(V,C),\qquad
 \lambda\otimes v\longmapsto
 \bigl(\varphi\longmapsto\lambda(\varphi(v))\bigr).
\]
Apply this isomorphism termwise to the resolution $R$ of
Lemma~\ref{res:base}. Exactness of vector-space duality shows that
$L\otimes_C R$ has homology only in degree two, where it is
$D(s^{\mathrm r})=s$. Moreover, $L$ acts by zero on the entire complex,
since $L^2=0$. Good truncations therefore give a $T$-linear derived
identification
\[
 L\otimes_C R\simeq s[2].
\]

Regard $0\to L\to T\to C\to0$ as a sequence of $(T,C)$-bimodules.
Tensoring it termwise with $R$ gives a short exact sequence of complexes
of left $T$-modules,
\[
 0\longrightarrow L\otimes_C R
 \longrightarrow T\otimes_C R
 \longrightarrow R\longrightarrow0,
\]
where the last complex is inflated along $T\to C$. Exactness uses
projectivity of each left $C$-module $R_i$; it does not require flatness
of $L$ or $T$ as right $C$-modules. Put
$\mathcal I=T\otimes_C R$. We obtain the derived triangle
\begin{equation}\label{res:triangle}
 s[2]\longrightarrow\mathcal I\longrightarrow s
 \xrightarrow{\,\delta\,}s[3].
\end{equation}
The middle complex has homology $s$ in degrees zero and two and no other
homology. Its terms are finite dimensional, so good truncation makes
it an object of the bounded derived category of finite $T$-modules.

Each term of $\mathcal I$ is a finite projective $T$-module. Because the
complex is bounded below homologically, it computes derived maps out of
it: a map from it to an acyclic complex is nullhomotopic by successive
projective lifts starting in its bottom degree. Termwise induction and
restriction adjunction thus gives, for every integer $i$,
\[
 B^i:=\Hom_{D^b(T\text{-mod})}(\mathcal I,s[i])
 =H^i\Hom_C(R,s)
 =\begin{cases}\kk&i=0,\\0&i\ne0.\end{cases}
\]

Set $E^i=\Hom_{D^b(T\text{-mod})}(s,s[i])$. These are ordinary Ext
groups for $i\geq0$ and vanish for $i<0$. Applying
$\Hom(-,s[i])$ to \eqref{res:triangle} gives the exact segment
\begin{equation}\label{res:recurrence}
 B^{i-1}\longrightarrow E^{i-3}
 \xrightarrow{\,\delta^*\,}E^i\longrightarrow B^i.
\end{equation}
For $i=1$, the groups $E^{-2}$ and $B^1$ vanish, so $E^1=0$.
For every $i\geq2$, both displayed $B$-groups vanish; hence
$\delta^*:E^{i-3}\to E^i$ is an isomorphism. This includes $E^2=0$,
and at $i=3$ it sends $\id_s$ to the nonzero class $\delta$.

More precisely, if $\beta:s\to s[i-3]$, its image under this connecting
map is
\[
 \delta^*(\beta)=\beta[3]\circ\delta.
\]
Thus \eqref{res:recurrence} is multiplication by $\delta$ in the Yoneda
algebra. Every power $\delta^m$ is nonzero and spans $E^{3m}$, and all
other nonnegative degrees vanish. We have proved
$\Ext_T^*(s,s)=\kk[\delta]$ as an algebra.

Lemma~\ref{coc:data} shows that $p(s)$ is nonzero in degree three: the
simple-valued cochain evaluates to $q^3\ne0$ on the cycle
$q^2[t|x|J]+[t|y|J]$. Hence $p(s)$ is a nonzero scalar multiple of
$\delta$ and is itself a polynomial generator. This proves the stated
normalization $\tau=p(s)$ without choosing an isomorphism
$L\otimes_C R\simeq s[2]$ with a prescribed scalar.
\end{proof}

\subsection{Two independent degree-three generators}

We now pass to the symmetric algebra and simple used in the remaining
construction:
\[
 A=T\otimes_{\kk}T,\qquad X=s\otimes_{\kk}s.
\]
The tensor product of the symmetrizing forms makes $A$ symmetric, and
$\dim_{\kk}A=400$. Let $p_1,p_2$ be the external copies of the
bimodule class $p$ on the two factors, and write
$\tau_i=p_i(X)$ for their evaluations.

\begin{lemma}\label{res:tensor}
The simple left $A$-module $X$ is nonprojective, and
\[
 \bigoplus_{a\geq0}\stHom{A}{X}{X[a]}
 =\kk[\tau_1,\tau_2],\qquad |\tau_1|=|\tau_2|=3.
\]
The multiplication is composition with shifts, and its positive-degree
part agrees with the ordinary Ext algebra.
\end{lemma}

\begin{proof}
Tensor two projective $T$-resolutions of $s$ over $\kk$. In each total
degree the direct sum is finite, so this is a resolution of $X$ by
finite projective $A$-modules. The Hom complex is the total tensor
product of the two Hom complexes over $T$. Splitting complexes of
vector spaces gives the K\"unneth isomorphism. Tensoring the comparison
maps identifies multiplication as well; the external maps commute in
characteristic two. Lemma~\ref{res:polynomial} therefore gives
\[
 \Ext_A^*(X,X)=\kk[\tau_1,\tau_2].
\]
In particular $X$ has a nonzero positive self-extension and cannot be
projective. Its ordinary endomorphism ring is $\kk$, since $X$ is one
dimensional. Its identity cannot factor through a projective module:
such a factorization would make $X$ a direct summand of that projective.
Thus its stable endomorphism ring is also $\kk$. Positive stable Hom
agrees with Ext over the symmetric algebra $A$, giving the result.
\end{proof}

We have obtained two polynomial classes acting on $X$, together with
specific bimodule representatives for them. The next step uses their
cones to leave only two stable Hom groups; their degrees will determine
the shift in the subsequent fiber construction.
